\section{Per-Iteration Periodicity Checking under High-Precision Arithmetic on the GPU}
\label{sec:hp-periodicity}

The \FractalShark{} fused v2 GPU reference-orbit implementation
(\cref{subsec:ref-orbit-gpu}) evaluates the Mandelbrot recurrence
\(
z_{n+1} = z_n^2 + c
\)
using high-precision arithmetic for the state \(z_n\), while interleaving a
reduced-precision periodicity/termination test before the high-precision
recurrence stages.  Early termination can avoid later arithmetic, but the
check and its required cooperative-grid synchronization have a measurable cost
that depends on precision and launch configuration.

Application dispatch to the CUDA engine is manual.  Its stored-output type \code{T} may be
\code{HDRFloat<float>}, \code{HDRFloat<double>}, or
\code{HDRFloat<CudaDblflt<dblflt>>}.  The CUDA path forces
\code{PExtras::Disable} and \code{ReuseMode::DontSaveForReuse}; neither produces
a compressed or intermediate-reuse orbit.  These are policy constraints of the
CUDA application path.

\paragraph{State extraction and logging.}

At iteration \(n\), the high-precision state \((x_n, y_n)\) is available in the
persistent kernel's reference state. The periodicity checker converts
these values into a reduced high-dynamic-range scalar format,
\(\mathrm{HDR}(\cdot)\), via a lossy but monotone extraction
(\code{ToHDRFloat}) and stores the resulting pair in a per-iteration trace
buffer:

\begin{equation}
  \widehat{x}_n \leftarrow \mathrm{HDR}(x_n), \quad
  \widehat{y}_n \leftarrow \mathrm{HDR}(y_n), \quad
  \code{OutputIters}[n] \leftarrow (\widehat{x}_n, \widehat{y}_n).
  \label{eq:hdr-extraction}
\end{equation}

This trace is used both for downstream diagnostics and to provide a compact,
device-side reduced representation of the orbit prefix.

\paragraph{Derivative-based periodicity criterion.}

In addition to the orbit state, the computation propagates the complex derivative
\(
\frac{dz}{dc}
\)
along the orbit (cf.\ the CPU-side counterpart in \cref{subsec:ref-orbit-periodicity}). For the Mandelbrot map \(f(z)=z^2+c\), the chain rule yields

\begin{equation}
  \frac{d z_{n+1}}{d c} = f'(z_n)\frac{d z_n}{d c} + 1
  = 2 z_n \frac{d z_n}{d c} + 1,
  \label{eq:dzdc-recurrence}
\end{equation}

with initialization \(\frac{dz_0}{dc}=0\) when \(z_0=0\).
Writing \(z_n = x_n + i y_n\) and introducing the shorthand
\(d_n \equiv \frac{\partial z_n}{\partial c}\) with real and imaginary parts
\(d_{x,n}\) and \(d_{y,n}\), expanding
\eqref{eq:dzdc-recurrence} into real and imaginary parts gives:

\begin{align}
  d_{x,n+1} &= 2\,(x_n\, d_{x,n} - y_n\, d_{y,n}) + 1, \label{eq:dzdc-real}\\
  d_{y,n+1} &= 2\,(x_n\, d_{y,n} + y_n\, d_{x,n}). \label{eq:dzdc-imag}
\end{align}

In the CUDA engine, the orbit components are \code{HpSharkFloat}, but the
application periodicity state \code{dzdcX}, \code{dzdcY}, and \code{RadiusY}
uses \code{HDRFloat}.  The periodicity decision is performed on reduced HDR
magnitudes:

\begin{equation}
  \lVert z_n\rVert_\infty \approx \max(|\widehat{x}_n|, |\widehat{y}_n|), \qquad
  \left\lVert\frac{dz_n}{dc}\right\rVert_\infty \approx \max(|\widehat{d x}_n|, |\widehat{d y}_n|),
  \label{eq:reduced-norms}
\end{equation}

where hats denote reduced HDR representations after a normalization/reduction
step (\code{HdrReduce}) to ensure stable comparisons.

The checker uses a sufficient condition of the form

\begin{equation}
  \lVert z_n\rVert_\infty < 2\,R\,\left\lVert\frac{dz_n}{dc}\right\rVert_\infty,
  \label{eq:periodicity-ineq}
\end{equation}

where \(R\) is a problem-dependent bound (stored as \code{RadiusY} in the
reference structure). Intuitively, \eqref{eq:periodicity-ineq} compares the
current orbit magnitude against a scaled sensitivity radius implied by the
derivative growth: when the orbit state becomes small relative to the local
linearization bound, further iteration is treated as converging toward a
previous state and the algorithm flags a periodic cycle.
Concretely, the device computes

\begin{equation}
  n_2 \leftarrow \max(|\widehat{x}_n|, |\widehat{y}_n|), \quad
  r_0 \leftarrow \max(|\widehat{d x}_n|, |\widehat{d y}_n|), \quad
  n_3 \leftarrow 2\,R\,r_0,
  \label{eq:gpu-periodicity-vars}
\end{equation}

and terminates when \(n_2 < n_3\), setting \code{PeriodicityStatus =
PeriodFound} and recording \code{OutputIterCount = n+1}.

\paragraph{Escape test.}
We additionally apply a conventional escape-radius test in the same reduced
domain. After forming the next-state candidates
\(
\widehat{x}_n + \widehat{c}_x
\)
and
\(
\widehat{y}_n + \widehat{c}_y
\),
the checker evaluates
\begin{equation}
  \lVert\widehat{z}_n + \widehat{c}\rVert_2^2
  = (\widehat{x}_n + \widehat{c}_x)^2 + (\widehat{y}_n + \widehat{c}_y)^2,
  \label{eq:gpu-escape-norm}
\end{equation}
and declares \code{Escaped} when this value exceeds a fixed threshold
(\code{256.0} in the current implementation). As with the periodicity
criterion, all comparisons are performed on reduced positive HDR values to avoid
high-precision control-flow costs on the device.

\paragraph{Control-flow placement and synchronization.}
Periodicity checking is invoked once per iteration at the top of the reference
step, before the fused high-precision stages compute
\(z_{n+1}\). In the production kernel, a single distinguished thread
(\code{block 0, thread 0}) executes the checker and updates the shared
termination flag stored in the global \code{reference} structure. The kernel
then executes a cooperative-grid barrier:
\begin{equation}
  \code{grid.sync()},
  \label{eq:grid-sync}
\end{equation}
ensuring that (i) all blocks observe a consistent \code{PeriodicityStatus} and
(ii) any updates to \code{dzdcX}, \code{dzdcY}, and the trace buffer are
visible before subsequent work proceeds. If termination is requested, the
iteration loop exits immediately; otherwise, the kernel proceeds into the
fused NTT recurrence that produces the next high-precision state.

This organization isolates termination logic from the fused recurrence pipeline
while maintaining grid-wide early-exit behavior.  The checker publishes its
status through one explicit \code{grid.sync()} per iteration.  The recurrence uses the same publication
boundary to expose the iteration plan built by the leader.  The synchronization
cost is amortized by the subsequent high-precision work but is not assumed to be
negligible.

\paragraph{Alternative host-side periodicity checking.}
An alternative design would be to retain the authoritative high-precision orbit
state entirely on the device, periodically copy the full-precision values
\((x_n, y_n, dz_n/dc)\) back to the host, and perform periodicity and escape
testing on the CPU using existing MPIR-based logic. While conceptually simpler,
this approach is impractical at large precisions: each iteration would require
transferring hundreds to thousands of limbs per component over the PCIe bus,
and the cumulative bandwidth and synchronization costs would dominate the total
runtime. Moreover, host-side checking would introduce additional latency and
force a tighter coupling between device execution and CPU control flow. By
performing periodicity detection directly on the GPU using reduced HDR
representations, the implementation avoids repeated large data transfers while
preserving an early-termination mechanism that is both inexpensive and tightly
integrated with the cooperative-grid execution model.

The reference-orbit algorithms described above rely on high-precision arithmetic
that must execute efficiently on the GPU.  The next chapter details the
NTT-based arithmetic engine that provides this capability.
